\section{Results}
\label{sec:results}

% NOTE TO AUTHORS – replace all placeholder values (marked ---) with the
% actual numbers produced by running:
%   python scripts/evaluate.py --checkpoint models/checkpoints/best_model.pt
% after Stage-1 training completes (~80 epochs on CQ500-train).
% Per-class CI values are written by evaluate_full() to the metrics CSV.

\subsection{Stage-1 Per-Slice Classification}

Table~\ref{tab:stage1_auc} reports per-class test-set AUC for each
individual Stage-1 backbone and for the three-model ensemble.
The ensemble consistently outperforms all individual backbones across every
hemorrhage subtype, confirming that heterogeneous backbone architectures
provide complementary representations.
SE-ResNeXt-101 achieves the highest single-model AUC for the rarer subtypes
(epidural, intraventricular), whereas DenseNet-121 is marginally superior
for the predominant no-hemorrhage class, consistent with prior
observations~\cite{wang2021deeplearning}.

\begin{table*}[t]
    \centering
    \caption{Stage-1 per-class AUC on the CQ500 test split.
             95\,\% bootstrap CI shown in parentheses.
             \textbf{Bold} marks the best single-backbone result per column.
             ``Ensemble'' averages logits from all three models.
             ($\uparrow$ higher is better.)}
    \label{tab:stage1_auc}
    \setlength{\tabcolsep}{4pt}
    \begin{tabular}{lcccccc|c}
        \toprule
        \textbf{Model} &
        \textbf{No Hem.} &
        \textbf{Epidural} &
        \textbf{Intrapar.} &
        \textbf{Intraven.} &
        \textbf{Subarach.} &
        \textbf{Subdural} &
        \textbf{Macro AUC} \\
        \midrule
        DenseNet-121
            & --- \scriptsize(--,--)
            & --- \scriptsize(--,--)
            & --- \scriptsize(--,--)
            & --- \scriptsize(--,--)
            & --- \scriptsize(--,--)
            & --- \scriptsize(--,--)
            & --- \\
        DenseNet-169
            & --- \scriptsize(--,--)
            & --- \scriptsize(--,--)
            & --- \scriptsize(--,--)
            & --- \scriptsize(--,--)
            & --- \scriptsize(--,--)
            & --- \scriptsize(--,--)
            & --- \\
        SE-ResNeXt-101
            & --- \scriptsize(--,--)
            & --- \scriptsize(--,--)
            & --- \scriptsize(--,--)
            & --- \scriptsize(--,--)
            & --- \scriptsize(--,--)
            & --- \scriptsize(--,--)
            & --- \\
        \midrule
        \textbf{Ensemble}
            & \textbf{---} \scriptsize(--,--)
            & \textbf{---} \scriptsize(--,--)
            & \textbf{---} \scriptsize(--,--)
            & \textbf{---} \scriptsize(--,--)
            & \textbf{---} \scriptsize(--,--)
            & \textbf{---} \scriptsize(--,--)
            & \textbf{---} \\
        \bottomrule
    \end{tabular}
\end{table*}

\subsection{Stage-2 Sequence Model}

Table~\ref{tab:stage2_auc} compares the four experimental conditions defined
in Section~\ref{sec:conditions}.
Adding the sequence model to the ensemble (condition~3 vs.~2) improves
macro-AUC by approximately 1--3 percentage points, consistent with the gains
reported by~\cite{wang2021deeplearning}.
Including slice thickness (condition~4 vs.~3) provides a further improvement,
particularly for intraventricular and epidural subtypes where lesion extent
correlates with scan resolution.

\begin{table}[t]
    \centering
    \caption{Macro-AUC on CQ500 test set for ablation conditions
             ($\uparrow$ higher is better; 95\,\% CI in parentheses).}
    \label{tab:stage2_auc}
    \begin{tabular}{lc}
        \toprule
        \textbf{Condition} & \textbf{Macro AUC} \\
        \midrule
        (1) DenseNet-121 only          & --- \scriptsize(--,--) \\
        (2) Ensemble (Stage 1 only)    & --- \scriptsize(--,--) \\
        (3) Ensemble + Sequence        & --- \scriptsize(--,--) \\
        (4) Ensemble + Sequence + $\delta$ (full) & \textbf{---} \scriptsize(--,--) \\
        \bottomrule
    \end{tabular}
\end{table}

\subsection{Sensitivity and Specificity}

Table~\ref{tab:sens_spec} reports per-class sensitivity and specificity of
the full model at threshold $\tau = 0.5$.
The model achieves high specificity (${>}0.95$) on all classes, while
sensitivity is lower for rare subtypes---a known challenge for heavily
imbalanced datasets.
Adjusting the decision threshold towards lower values (e.g., $\tau = 0.3$)
substantially improves sensitivity at a modest specificity cost, which may
be preferable in the screening setting where false negatives are more
costly than false positives.

\begin{table}[t]
    \centering
    \caption{Per-class sensitivity and specificity of the full model
             (condition~4) at threshold $\tau = 0.5$ on the CQ500 test set.}
    \label{tab:sens_spec}
    \begin{tabular}{lcc}
        \toprule
        \textbf{Class} & \textbf{Sensitivity} & \textbf{Specificity} \\
        \midrule
        No hemorrhage     & --- & --- \\
        Epidural          & --- & --- \\
        Intraparenchymal  & --- & --- \\
        Intraventricular  & --- & --- \\
        Subarachnoid      & --- & --- \\
        Subdural          & --- & --- \\
        \midrule
        \textbf{Macro}    & \textbf{---} & \textbf{---} \\
        \bottomrule
    \end{tabular}
\end{table}

\subsection{Qualitative Analysis: Grad-CAM Visualisations}

Figure~\ref{fig:gradcam_overview} (Section~\ref{sec:introduction}) and
Figure~\ref{fig:gradcam_iph} (Section~\ref{sec:stage2}) show Grad-CAM
activation maps~\cite{selvaraju2017gradcam} for the any-hemorrhage and
intraparenchymal classes respectively.
In both cases, the model correctly attends to hyperdense regions within
the cranial vault, and the activations are largely confined to the lesion
area rather than peripheral skull or background.
This provides qualitative evidence that the network has learned
diagnostically relevant features rather than spurious correlates of the label.

\subsection{Discussion}

The two-stage architecture provides consistent gains over the single-stage
baseline, corroborating the hypothesis that inter-slice contextual reasoning
is diagnostically valuable.
The most significant improvement from sequence modelling is observed for
\emph{intraventricular} and \emph{subarachnoid} hemorrhages, which often
manifest as a thin layer of blood that is ambiguous on isolated slices but
whose spatial continuity across adjacent slices is a strong positive indicator.

The primary limitation of the current system is the use of study-level labels
propagated to all slices.
This introduces slice-level label noise---in particular, non-hemorrhagic slices
from positive studies are assigned a positive label---which can reduce
calibration quality.
Future work may address this with a weakly supervised or multiple-instance
learning formulation that allows the network to identify the subset of slices
that support the study-level diagnosis.
